%% LyX 2.4.4 created this file.  For more info, see https://www.lyx.org/.
%% Do not edit unless you really know what you are doing.
\documentclass[english]{article}
\usepackage[T1]{fontenc}
\usepackage[latin9]{inputenc}
\usepackage{units}
\usepackage{enumitem}
\usepackage{amsmath}
\usepackage{amsthm}
\usepackage{cancel}
\usepackage{geometry}
\geometry{verbose,tmargin=1in,bmargin=1in,lmargin=1in,rmargin=1in}

\makeatletter
%%%%%%%%%%%%%%%%%%%%%%%%%%%%%% Textclass specific LaTeX commands.
\numberwithin{equation}{section}
\numberwithin{figure}{section}
\newlength{\lyxlabelwidth}      % auxiliary length 

%%%%%%%%%%%%%%%%%%%%%%%%%%%%%% User specified LaTeX commands.
\usepackage{fancyhdr}
\usepackage{lastpage}
\pagestyle{fancy}
\usepackage{enumitem}
\usepackage{ifthen}
%sets math fonts to be more readable
\everymath=\expandafter{\the\everymath\displaystyle}
%\DeclareMathSizes{10}{10}{10}{10}
\usepackage{multicol}

\usepackage{tikz}
\usetikzlibrary{plotmarks}
\usepackage{pgfplots}
\pgfplotsset{compat=newest}

\usepackage{calc}
\usepackage[nomessages]{fp}

\usepackage{advdate}    % Advancing/saving dates
\usepackage{datetime}   % Dates formatting
\usepackage{datenumber} % Counters for dates
% Added by lyx2lyx
\setlength{\parskip}{\smallskipamount}
\setlength{\parindent}{0pt}

\makeatother

\usepackage{babel}
\begin{document}
\renewcommand{\author}{Dr.\,James Ashe}

\newcommand{\classNum}{137}

\newcommand{\classSec}{B}

\newcommand{\nameType}{Name:}

\newcommand{\examType}{worksheet 4.1-4.3}

\renewcommand{\date}{\formatdate{22}{4}{2013}}

%%First page's header%%%%%%%%%%%%%%%%%%%%%%
\fancypagestyle{firststyle} %%header settings for first pagr
{
	\fancyhead[L]{MTH \classNum{}(\classSec{}) \author{}\\
	\textbf{\examType{}} \date{}}
	\fancyhead[C]{\textbf{\nameType}}
	\setlength{\headsep}{14pt}
}
%%%%%%%%%%%%%%%%%%%%%%%%%%%%%%%%%%%%%%%%%%

%%Next page's header%%%%%%%%%%%%%%%%%%%%%%%
\setlist{nolistsep}
\lhead{\classNum{}(\classSec{})} 
\chead{\textbf{\nameType}} 
\cfoot{\thepage{}~of~\pageref{LastPage}} 
\setlength{\headsep}{5pt} 
%%%%%%%%%%%%%%%%%%%%%%%%%%%%%%%%%%%%%%%%%%%

%%Various settings; see the preamble for other settings%%%%%
%%\setenumerate[0]{leftmargin=12pt,itemindent=0pt,labelwidth=0pt}
\renewcommand{\labelenumi}{\textbf{\arabic{enumi}.}}
\renewcommand{\labelenumii}{\textbf{\alph{enumii}.}}
\thispagestyle{firststyle}
%%%%%%%%%%%%%%%%%%%%%%%%%%%%%%%%%%%%%%%%%%
\newcommand{\xmax}{10}
\newcommand{\xmin}{-10}
\newcommand{\xstep}{0} %must be xmin+n
\newcommand{\ymax}{10}
\newcommand{\ymin}{-10}
\newcommand{\ystep}{0} %must be ymin+n

\newcommand{\Dspace}{\vspace{1cm}}
\begin{enumerate}
\item Evaluate the following:

\begin{enumerate}
\item $-2^{-3}$\vfill
\item $\left(-\frac{2}{3}\right)^{-2}$\vfill
\item $8^{-\nicefrac{4}{3}}$\vfill
\item $f(x)=5^{2-3x}$, find $f(1)$.\vfill
\item ${\displaystyle \lim_{x\rightarrow\infty}3^{x}}$\vfill
\item ${\displaystyle \lim_{x\rightarrow-\infty}3^{x}}$\vfill
\item ${\displaystyle \lim_{x\rightarrow\infty}3^{-x}}$\vfill
\end{enumerate}
\end{enumerate}
\begin{enumerate}[resume]
\item Graph the following:\begin{multicols}{2}

\begin{enumerate}
\item $f(x)=3^{-x}$\\
\renewcommand{\xmax}{10}
\renewcommand{\xmin}{-10}
\renewcommand{\xstep}{0} %must be xmin+n
\renewcommand{\ymax}{10}
\renewcommand{\ymin}{-10}
\renewcommand{\ystep}{0} %must be ymin+n
%%%%%%%
\begin{tikzpicture} 
\begin{axis}[height=5.5cm, 
	width=5.5cm, 
	axis lines=middle,
	minor x tick num={4},
	minor y tick num={4},
	grid=both,
	%xtick={0,50,...,250},
	%ytick={0,10,20,...,40},
	%axis line style={-latex}, 
	xmin=\xmin,xmax=\xmax, 
	ymin=\ymin,ymax=\ymax,
	%restrict y to domain=-1:10,
	no markers, 
	xlabel={$x$},ylabel={$y$}, 
	%every axis x label/.append style={anchor=north}, 
	%every axis y label/.append style={anchor=east}
	]
	
	%%\addplot[color=red,solid,mark=*] coordinates {(-3,-4)};
	%%\addplot[domain=\xmin:\xmax,thick,samples=100,draw=red,<->]{-(x+2)*(x-2)*(x-5)} node[pos=.65,right] {$f(x)$};
	
\end{axis}
\end{tikzpicture}

\begin{enumerate}
\item What is the domain and range of $f(x)$?\columnbreak
\end{enumerate}
\item $g(x)=3^{-(x-1)}+2$\\
\renewcommand{\xmax}{10}
\renewcommand{\xmin}{-10}
\renewcommand{\xstep}{0} %must be xmin+n
\renewcommand{\ymax}{10}
\renewcommand{\ymin}{-10}
\renewcommand{\ystep}{0} %must be ymin+n
%%%%%%%
\begin{tikzpicture} 
\begin{axis}[height=5.5cm, 
	width=5.5cm, 
	axis lines=middle,
	minor x tick num={4},
	minor y tick num={4},
	grid=both,
	%xtick={0,50,...,250},
	%ytick={0,10,20,...,40},
	%axis line style={-latex}, 
	xmin=\xmin,xmax=\xmax, 
	ymin=\ymin,ymax=\ymax,
	%restrict y to domain=-1:10,
	no markers, 
	xlabel={$x$},ylabel={$y$}, 
	%every axis x label/.append style={anchor=north}, 
	%every axis y label/.append style={anchor=east}
	]
	
	%%\addplot[color=red,solid,mark=*] coordinates {(-3,-4)};
	%%\addplot[domain=\xmin:\xmax,thick,samples=100,draw=red,<->]{-(x+2)*(x-2)*(x-5)} node[pos=.65,right] {$f(x)$};
	
\end{axis}
\end{tikzpicture}

\begin{enumerate}
\item What is the domain and range of $g(x)?$
\end{enumerate}
\end{enumerate}
\end{multicols}
\item Graph the following:\begin{multicols}{2}

\begin{enumerate}
\item $f(x)=\log_{3}x$\\
\renewcommand{\xmax}{10}
\renewcommand{\xmin}{-10}
\renewcommand{\xstep}{0} %must be xmin+n
\renewcommand{\ymax}{10}
\renewcommand{\ymin}{-10}
\renewcommand{\ystep}{0} %must be ymin+n
%%%%%%%
\begin{tikzpicture} 
\begin{axis}[height=5.5cm, 
	width=5.5cm, 
	axis lines=middle,
	minor x tick num={4},
	minor y tick num={4},
	grid=both,
	%xtick={0,50,...,250},
	%ytick={0,10,20,...,40},
	%axis line style={-latex}, 
	xmin=\xmin,xmax=\xmax, 
	ymin=\ymin,ymax=\ymax,
	%restrict y to domain=-1:10,
	no markers, 
	xlabel={$x$},ylabel={$y$}, 
	%every axis x label/.append style={anchor=north}, 
	%every axis y label/.append style={anchor=east}
	]
	
	%%\addplot[color=red,solid,mark=*] coordinates {(-3,-4)};
	%%\addplot[domain=\xmin:\xmax,thick,samples=100,draw=red,<->]{-(x+2)*(x-2)*(x-5)} node[pos=.65,right] {$f(x)$};
	
\end{axis}
\end{tikzpicture}

\begin{enumerate}
\item What is the domain and range of $f(x)$?\columnbreak
\end{enumerate}
\item $g(x)=\log_{3}\left(x-1\right)$\\
\renewcommand{\xmax}{10}
\renewcommand{\xmin}{-10}
\renewcommand{\xstep}{0} %must be xmin+n
\renewcommand{\ymax}{10}
\renewcommand{\ymin}{-10}
\renewcommand{\ystep}{0} %must be ymin+n
%%%%%%%
\begin{tikzpicture} 
\begin{axis}[height=5.5cm, 
	width=5.5cm, 
	axis lines=middle,
	minor x tick num={4},
	minor y tick num={4},
	grid=both,
	%xtick={0,50,...,250},
	%ytick={0,10,20,...,40},
	%axis line style={-latex}, 
	xmin=\xmin,xmax=\xmax, 
	ymin=\ymin,ymax=\ymax,
	%restrict y to domain=-1:10,
	no markers, 
	xlabel={$x$},ylabel={$y$}, 
	%every axis x label/.append style={anchor=north}, 
	%every axis y label/.append style={anchor=east}
	]
	
	%%\addplot[color=red,solid,mark=*] coordinates {(-3,-4)};
	%%\addplot[domain=\xmin:\xmax,thick,samples=100,draw=red,<->]{-(x+2)*(x-2)*(x-5)} node[pos=.65,right] {$f(x)$};
	
\end{axis}
\end{tikzpicture}

\begin{enumerate}
\item What is the domain and range of $g(x)?$
\end{enumerate}
\end{enumerate}
\end{multicols}\newpage
\item Evaluate the following:

\begin{enumerate}
\item $\left(\frac{1}{6}\right)^{x}=36$\vfill
\item $\log_{3}\left(27\right)$\vfill
\item $\log_{\sqrt{10}}1$\vfill
\item $\ln\left(\frac{1}{e^{2}}\right)$\vfill\Dspace
\item Write $z=\log_{6}5$ as an exponential equation.\vfill
\item ${\displaystyle \lim_{x\rightarrow0^{+}}\log_{3}x}$\vfill
\end{enumerate}
\item If \$1000 is invested in an account at $5\%$ annual interest compounded
monthly, what is the amount in the account after 6 years?\vfill\Dspace\Dspace\Dspace\newpage
\item Solve for $x$. Write your answer as an exact value and approximate
value.

\begin{enumerate}
\item $5^{x}=2$\vfill
\item $\left(3.5\right)^{2x}=10$\vfill
\item $2\ln x=7$\vfill
\item $\log_{2}x=\sqrt{5}$\vfill
\item $\left(1+x\right)^{4}=3.1$ (only approximate)\vfill
\end{enumerate}
\item If \$1000 is invested in an account at 5\% compounded continuously,
how long will it take to double the investment?\vfill\Dspace\Dspace\newpage
\item Use the properties of logarithms to expand the expression as much
as possible.

\begin{enumerate}[resume]
\item $\log_{2}\left(8x\right)$\vfill
\item $\ln\left(\frac{e^{6}}{11}\right)$\vfill
\item $\ln\left(\frac{\sqrt[3]{xy}}{t^{2}}\right)$\vfill
\end{enumerate}
\item Use the properties of logarithms to rewrite the expression as a single
logarithm.

\begin{enumerate}[resume]
\item $\log_{2}\left(5\right)+3\log_{2}\left(x\right)$\vfill
\item $\frac{1}{3}\ln(6)-\frac{1}{3}\ln\left(2\right)$\vfill
\item $\log\left(2\right)+\log\left(3\right)+\log\left(y\right)-\log\left(z\right)$\vfill
\item $\frac{1}{2}\left[\log\left(x\right)+\log\left(y\right)\right]-\log\left(z\right)$\vfill
\end{enumerate}
\end{enumerate}

\end{document}
